\chapter{2018年北京卷（文）}

\section{单选题}

\begin{problem}
已知集合 \(A=\{x\mid \lvert x\rvert<2\}\)，\(B=\{-2,0,1,2\}\)，则 \(A\cap B=\)（\quad）.

\choices
    {\(\{0,1\}\)}
    {\(\{-1,0,1\}\)}
    {\(\{-2,0,1,2\}\)}
    {\(\{-1,0,1,2\}\)}
\end{problem}

\begin{answer}
A
\end{answer}

\begin{solution}
由 \(\lvert x\rvert<2\) 得 \(-2<x<2\)，所以 \(A\) 中包含 \(0\) 和 \(1\)，但不包含 \(-2\) 和 \(2\). 因此 \(A\cap B=\{0,1\}\)，故选 A.
\end{solution}

\begin{problem}
在复平面内，复数 \(\dfrac1{1-\i}\) 的共轭复数对应的点位于（\quad）.

\choices
    {第一象限}
    {第二象限}
    {第三象限}
    {第四象限}
\end{problem}

\begin{answer}
D
\end{answer}

\begin{solution}
\(\displaystyle \frac{1}{1-\i}=\frac{1+\i}{(1-\i)(1+\i)}=\frac{1+\i}{2}\)，其共轭复数为 \(\displaystyle \frac{1-\i}{2}\). 对应点的坐标为 \(\left(\dfrac12,-\dfrac12\right)\)，实部为正、虚部为负，位于第四象限，故选 D.
\end{solution}


\begin{problem}
执行如图所示的程序框图，输出的 \(s\) 值为（\quad）.

\begingroup
\bitmapfigure[width=0.32\linewidth]{img/2018/beijing_liberal/q3_fig1.png}
\endgroup

\choices
    {\(\dfrac12\)}
    {\(\dfrac56\)}
    {\(\dfrac76\)}
    {\(\dfrac7{12}\)}
\end{problem}

\begin{answer}
B
\end{answer}

\begin{solution}
程序开始时 \(k=1\)，\(s=1\). 第一次循环将 \(k=1\) 代入，得 \(s=1+(-1)^1\cdot\dfrac1{1+1}=\dfrac12\)，随后 \(k=2\)，因 \(k<3\) 继续循环. 第二次循环得 \(s=\dfrac12+(-1)^2\cdot\dfrac1{1+2}=\dfrac56\)，随后 \(k=3\)，满足 \(k\geq3\)，输出 \(s=\dfrac56\)，故选 B.
\end{solution}


\begin{problem}
设 \(a\)，\(b\)，\(c\)，\(d\) 是非零实数，则“\(ad=bc\)”是“\(a\)，\(b\)，\(c\)，\(d\) 成等比数列”的（\quad）.

\choices
    {充分而不必要条件}
    {必要而不充分条件}
    {充分必要条件}
    {既不充分也不必要条件}
\end{problem}

\begin{answer}
B
\end{answer}

\begin{solution}
若 \(a\)，\(b\)，\(c\)，\(d\) 成等比数列，设公比为 \(q\)，则 \(b=aq\)，\(c=aq^2\)，\(d=aq^3\)，从而 \(ad=a^2q^3=bc\)，所以“\(a\)，\(b\)，\(c\)，\(d\) 成等比数列”是“\(ad=bc\)”的必要条件.

但反例 \(a=1\)，\(b=1\)，\(c=2\)，\(d=2\) 满足 \(ad=bc=2\)，而相邻两项的比依次为 \(1\)，\(2\)，\(1\)，不是等比数列. 因此该条件必要而不充分，故选 B.
\end{solution}


\begin{problem}
“十二平均律”是通用的音律体系，明代朱载堉最早用数学方法计算出半音比例，为这个理论的发展做出了重要贡献. 十二平均律将一个纯八度音程分成十二份，依次得到十三个单音，从第二个单音起，每一个单音的频率与它的前一个单音的频率的比都等于 \(\sqrt[12]{2}\). 若第一个单音的频率为 \(f\)，则第八个单音频率为（\quad）.

\choices
    {\(\sqrt[3]{2}\,f\)}
    {\(\sqrt[3]{2^2}\,f\)}
    {\(\sqrt[12]{2^5}\,f\)}
    {\(\sqrt[12]{2^7}\,f\)}
\end{problem}

\begin{answer}
D
\end{answer}

\begin{solution}
从第二个单音起，相邻两个单音频率的比为 \(\sqrt[12]{2}\)，所以第 \(n\) 个单音的频率为 \(f\left(\sqrt[12]{2}\right)^{n-1}\). 取 \(n=8\)，得第八个单音的频率为 \(f\left(\sqrt[12]{2}\right)^7=\sqrt[12]{2^7}\,f\)，故选 D.
\end{solution}


\begin{problem}
某四棱锥的三视图如图所示，在此四棱锥的侧面中，直角三角形的个数为（\quad）.

\begingroup
\bitmapfigure[float=inline,width=0.42\linewidth]{img/2018/beijing_liberal/q6_fig1.png}
\endgroup
\FloatBarrier

\choices
    {\(1\)}
    {\(2\)}
    {\(3\)}
    {\(4\)}
\end{problem}

\begin{answer}
C
\end{answer}

\begin{solution}
按三视图还原四棱锥，记顶点为 \(P\)，底面顶点按俯视图顺次为 \(A\)，\(B\)，\(C\)，\(D\). 由三视图可得 \(PA\perp\) 平面 \(ABCD\)，且 \(PA=2\)，\(AB=AD=2\)，\(BC=1\)，\(AC=CD=\sqrt5\).

因为 \(PA\perp\) 平面 \(ABCD\)，所以三角形 \(PAB\) 和三角形 \(PAD\) 都在点 \(A\) 处为直角. 又 \(PB^2=PA^2+AB^2=8\)，\(PC^2=PA^2+AC^2=9\)，且 \(BC^2=1\)，故 \(PB^2+BC^2=PC^2\)，三角形 \(PBC\) 在点 \(B\) 处为直角.

对于三角形 \(PCD\)，三边长分别为 \(PC=3\)，\(PD=\sqrt{PA^2+AD^2}=2\sqrt2\)，\(CD=\sqrt5\)，有 \(3^2\)，\((2\sqrt2)^2\)，\((\sqrt5)^2\) 中任意两数之和都不等于第三数，故它不是直角三角形. 因此侧面中直角三角形共有 \(3\) 个，故选 C.
\end{solution}


\begin{problem}
在平面直角坐标系中，\(\myarc{AB}\)，\(\myarc{CD}\)，\(\myarc{EF}\)，\(\myarc{GH}\) 是圆 \(x^2+y^2=1\) 上的四段弧（如图），点 \(P\) 在其中一段上，角 \(\alpha\) 以 \(Ox\) 为始边，\(OP\) 为终边，若 \(\tan\alpha<\cos\alpha<\sin\alpha\)，则 \(P\) 所在的圆弧是（\quad）.

\begingroup
\bitmapfigure[width=0.38\linewidth]{img/2018/beijing_liberal/q7_fig1.png}
\endgroup

\choices
    {\(\myarc{AB}\)}
    {\(\myarc{CD}\)}
    {\(\myarc{EF}\)}
    {\(\myarc{GH}\)}
\end{problem}

\begin{answer}
C
\end{answer}

\begin{solution}
设 \(P=(\cos\alpha,\sin\alpha)\). 若 \(P\) 在第一象限，则 \(0<\cos\alpha<1\)，而 \(\sin\alpha>\cos\alpha\) 时有 \(\tan\alpha=\dfrac{\sin\alpha}{\cos\alpha}>1>\cos\alpha\)，不可能满足 \(\tan\alpha<\cos\alpha\). 第四象限中 \(\cos\alpha>0>\sin\alpha\)，不满足 \(\cos\alpha<\sin\alpha\)；第三象限中 \(\tan\alpha>0>\cos\alpha\)，也不可能满足不等式.

因此只能有 \(\cos\alpha<0\)，\(\sin\alpha>0\)，即 \(P\) 在第二象限. 图中第二象限对应的圆弧为 \(\myarc{EF}\)，故选 C.
\end{solution}


\begin{problem}
设集合 \(A=\{(x,y)\mid x-y\ge1,\ ax+y>4,\ x-ay\le2\}\)，则（\quad）.

\choices
    {对任意实数 \(a\)，\((2,1)\in A\)}
    {对任意实数 \(a\)，\((2,1)\notin A\)}
    {当且仅当 \(a<0\) 时，\((2,1)\notin A\)}
    {当且仅当 \(a\le\dfrac32\) 时，\((2,1)\notin A\)}
\end{problem}

\begin{answer}
D
\end{answer}

\begin{solution}
将 \((x,y)=(2,1)\) 代入集合定义，第一条不等式为 \(2-1=1\geq1\)，恒成立；第二条不等式给出 \(2a+1>4\)，即 \(a>\dfrac32\)；第三条不等式给出 \(2-a\leq2\)，即 \(a\geq0\). 三者同时成立当且仅当 \(a>\dfrac32\).

所以 \((2,1)\notin A\) 当且仅当 \(a\leq\dfrac32\)，故选 D.
\end{solution}

\section{填空题}

\begin{problem}
设向量 \(\bs a=(1,0)\)，\(\bs b=(-1,m)\)，若 \(\bs a\perp(m\bs a-\bs b)\)，则 \(m=\fillinblank{}\).
\end{problem}

\begin{answer}
\(-1\)
\end{answer}

\begin{solution}
\(m\bs a-\bs b=m(1,0)-(-1,m)=(m+1,-m)\). 由垂直向量的数量积为零，得 \(\bs a\cdot(m\bs a-\bs b)=(1,0)\cdot(m+1,-m)=m+1=0\)，所以 \(m=-1\).
\end{solution}


\begin{problem}
已知直线 \(l\) 过点 \((1,0)\) 且垂直于 \(x\) 轴，若 \(l\) 被抛物线 \(y^2=4ax\) 截得的线段长度为 \(4\)，则抛物线的焦点坐标为\fillinblank{}.
\end{problem}

\begin{answer}
\((1,0)\)
\end{answer}

\begin{solution}
直线 \(l\) 垂直于 \(x\) 轴且过 \((1,0)\)，所以其方程为 \(x=1\). 代入抛物线方程得 \(y^2=4a\)，截得的弦长为两个交点纵坐标之差，即 \(2\sqrt{4a}=4\)，从而 \(a=1\). 抛物线 \(y^2=4ax\) 的焦点为 \((a,0)\)，故焦点坐标为 \((1,0)\).
\end{solution}


\begin{problem}
能说明“若 \(a>b\)，则 \(\dfrac1a<\dfrac1b\)”为假命题的一组 \(a\)，\(b\) 的值依次为\fillinblank{}.
\end{problem}

\begin{answer}
\(1\)，\(-1\)
\end{answer}

\begin{solution}
取 \(a=1\)，\(b=-1\)，则 \(a>b\)，但 \(\dfrac1a=1>-1=\dfrac1b\)，不满足题中结论 \(\dfrac1a<\dfrac1b\). 因此这组数说明原命题为假.
\end{solution}


\begin{problem}
若双曲线 \(\dfrac{x^2}{a^2}-\dfrac{y^2}{4}=1\)（\(a>0\)）的离心率为 \(\dfrac{\sqrt5}{2}\)，则 \(a=\fillinblank{}\).
\end{problem}

\begin{answer}
4
\end{answer}

\begin{solution}
双曲线的实半轴为 \(a\)，虚半轴为 \(2\)，所以焦半距满足 \(c^2=a^2+4\). 由离心率公式，\(\dfrac ca=\dfrac{\sqrt5}{2}\)，两边平方得 \(\dfrac{a^2+4}{a^2}=\dfrac54\)，即 \(4a^2+16=5a^2\)，所以 \(a^2=16\). 因为 \(a>0\)，故 \(a=4\).
\end{solution}


\begin{problem}
若 \(x\)，\(y\) 满足 \(x+1\le y\le2x\)，则 \(2y-x\) 的最小值是\fillinblank{}.
\end{problem}

\begin{answer}
3
\end{answer}

\begin{solution}
由 \(x+1\leq y\leq2x\) 可知 \(x+1\leq2x\)，所以 \(x\geq1\). 由于目标式 \(2y-x\) 随 \(y\) 增大而增大，对固定的 \(x\) 取最小允许值 \(y=x+1\)，于是 \(2y-x\geq2(x+1)-x=x+2\geq3\). 当 \(x=1\)，\(y=2\) 时等号成立，故最小值为 \(3\).
\end{solution}


\begin{problem}
若 \(\triangle ABC\) 的面积为 \(\dfrac{\sqrt3}{4}(a^2+c^2-b^2)\)，且 \(\angle C\) 为钝角，则 \(\angle B=\fillinblank{}\)；\(\dfrac ca\) 的取值范围是\fillinblank{}.
\end{problem}

\begin{answer}
\(60^\circ\)；\((2,+\infty)\)
\end{answer}

\begin{solution}
由余弦定理 \(a^2+c^2-b^2=2ac\cos B\)，又三角形面积 \(S=\dfrac12ac\sin B\)，故题设给出 \(\dfrac12ac\sin B=\dfrac{\sqrt3}{2}ac\cos B\). 因为 \(a\)，\(c>0\) 且 \(B\in(0,\pi)\)，所以 \(\sin B=\sqrt3\cos B\)，从而 \(B\) 为锐角且 \(\tan B=\sqrt3\)，得 \(B=60^\circ\).

又因 \(C\) 为钝角，且 \(A+B+C=180^\circ\)，所以 \(0^\circ<A<30^\circ\). 由正弦定理，\(\dfrac ca=\dfrac{\sin C}{\sin A}=\dfrac{\sin(120^\circ-A)}{\sin A}=\dfrac12+\dfrac{\sqrt3}{2}\cot A\). 由于 \(0^\circ<A<30^\circ\)，有 \(\cot A>\sqrt3\)，因此 \(\dfrac ca>2\). 当 \(A\to0^+\) 时 \(\dfrac ca\to+\infty\)，且 \(A\to30^-\) 时 \(\dfrac ca\to2^+\)，故取值范围为 \((2,+\infty)\).
\end{solution}

\section{解答题}

\begin{problem}

设 \(\{a_n\}\) 是等差数列，且 \(a_1=\ln2\)，\(a_2+a_3=5\ln2\).

\begin{enumerate}
\item 求 \(\{a_n\}\) 的通项公式；
\item 求 \(\e^{a_1}+\e^{a_2}+\cdots+\e^{a_n}\).
\end{enumerate}
\end{problem}

\begin{solution}
设等差数列 \(\{a_n\}\) 的公差为 \(d\). 由 \(a_2=a_1+d\)，\(a_3=a_1+2d\)，得
\[
a_2+a_3=2a_1+3d=5\ln2
\]
代入 \(a_1=\ln2\)，得到 \(d=\ln2\)，所以
\[
a_n=a_1+(n-1)d=\ln2+(n-1)\ln2=n\ln2
\]

又 \(\e^{a_k}=\e^{k\ln2}=2^k\)，因此
\[
\e^{a_1}+\e^{a_2}+\cdots+\e^{a_n}=2^1+2^2+\cdots+2^n=\frac{2(2^n-1)}{2-1}=2^{n+1}-2
\]
\end{solution}


\begin{problem}

已知函数 \(f(x)=\sin^2x+\sqrt3\sin x\cos x\).

\begin{enumerate}
\item 求 \(f(x)\) 的最小正周期；
\item 若 \(f(x)\) 在区间 \(\left[-\dfrac\pi3,m\right]\) 上的最大值为 \(\dfrac32\)，求 \(m\) 的最小值.
\end{enumerate}
\end{problem}

\begin{solution}
利用降幂公式和二倍角公式，得
\[
f(x)=\frac{1-\cos2x}{2}+\frac{\sqrt3}{2}\sin2x=\frac12+\sin\left(2x-\frac\pi6\right)
\]
正弦函数的自变量系数为 \(2\)，所以 \(f(x)\) 的最小正周期为 \(\dfrac{2\pi}{2}=\pi\).

因为 \(f(x)\leq\dfrac32\)，区间 \(\left[-\dfrac\pi3,m\right]\) 上的最大值等于 \(\dfrac32\) 的充要条件是该区间中存在 \(x\) 使 \(2x-\dfrac\pi6=\dfrac\pi2+2k\pi\)，其中 \(k\in\Z\). 令 \(t=2x-\dfrac\pi6\)，则当 \(x=-\dfrac\pi3\) 时 \(t=-\dfrac{5\pi}{6}\)，其右侧第一个使正弦值为 \(1\) 的点是 \(t=\dfrac\pi2\). 因此必须有 \(2m-\dfrac\pi6\geq\dfrac\pi2\)，即 \(m\geq\dfrac\pi3\). 取 \(m=\dfrac\pi3\) 时达到最大值，故 \(m\) 的最小值为 \(\dfrac\pi3\).
\end{solution}


\begin{problem}

电影公司随机收集了电影的有关数据，经分类整理得到下表：

\begin{center}
\begin{tabular}{c|cccccc}
电影类型 & 第一类 & 第二类 & 第三类 & 第四类 & 第五类 & 第六类 \\
\hline
电影部数 & 140 & 50 & 300 & 200 & 800 & 510 \\
好评率 & 0.4 & 0.2 & 0.15 & 0.25 & 0.2 & 0.1
\end{tabular}
\end{center}

好评率是指：一类电影中获得好评的部数与该类电影的部数的比值.

\begin{enumerate}
\item 从电影公司收集的电影中随机选取 \(1\) 部，求这部电影是获得好评的第四类电影的概率；
\item 随机选取 \(1\) 部电影，估计这部电影没有获得好评的概率；
\item 电影公司为增加投资回报，拟改变投资策略，这将导致不同类型电影的好评率发生变化. 假设表格中只有两类电影的好评率数据发生变化，那么哪类电影的好评率增加 \(0.1\)，哪类电影的好评率减少 \(0.1\)，使得获得好评的电影部数与样本中的电影总部数的比值达到最大？（只需写出结论）
\end{enumerate}
\end{problem}

\begin{solution}
样本中的电影总数为 \(140+50+300+200+800+510=2000\).

\begin{enumerate}
\item 第四类电影中获得好评的部数为 \(200\times0.25=50\)，所以所求概率为 \(\dfrac{50}{2000}=0.025\).
\item 样本中获得好评的电影部数为
\[
140\times0.4+50\times0.2+300\times0.15+200\times0.25+800\times0.2+510\times0.1=56+10+45+50+160+51=372
\]
所以没有获得好评的电影部数为 \(2000-372=1628\)，所求概率的估计值为 \(\dfrac{1628}{2000}=0.814\).
\item 若第 \(i\) 类好评率增加 \(0.1\)，第 \(j\) 类好评率减少 \(0.1\)，获得好评的电影部数的变化量为 \(0.1N_i-0.1N_j=0.1(N_i-N_j)\)，其中 \(N_i\)，\(N_j\) 分别为两类电影的部数. 要使变化量最大，应使增加好评率的类别部数最大、减少好评率的类别部数最小，即第五类增加 \(0.1\)，第二类减少 \(0.1\).
\end{enumerate}
\end{solution}


\begin{problem}

如图，在四棱锥 \(P-ABCD\) 中，底面 \(ABCD\) 为矩形，平面 \(PAD\perp\) 平面 \(ABCD\)，\(PA\perp PD\)，\(PA=PD\)，\(E\)，\(F\) 分别为 \(AD\)，\(PB\) 的中点.

\begin{enumerate}
\item 求证：\(PE\perp BC\)；
\item 求证：平面 \(PAB\perp\) 平面 \(PCD\)；
\item 求证：\(EF\parallel\) 平面 \(PCD\).
\end{enumerate}

\begingroup
\bitmapfigure[float=inline,width=0.62\linewidth]{img/2018/beijing_liberal/q18_fig1.png}
\endgroup
\FloatBarrier
\end{problem}

\begin{solution}
\begin{enumerate}
\item 因为 \(PA=PD\)，且 \(E\) 是 \(AD\) 的中点，所以在等腰三角形 \(PAD\) 中，\(PE\perp AD\). 又因为底面 \(ABCD\) 是矩形，所以 \(BC\parallel AD\)，从而 \(PE\perp BC\).
\item 过点 \(P\) 作直线 \(PG\parallel AB\). 因为 \(AB\parallel CD\)，所以 \(PG\parallel CD\)，且 \(PG\) 是平面 \(PAB\) 与平面 \(PCD\) 的交线. 底面平面与平面 \(PAD\) 垂直，且 \(AB\) 在底面内并垂直于交线 \(AD\)，由面面垂直的性质得 \(AB\perp\) 平面 \(PAD\)，所以 \(AB\perp PA\) 且 \(AB\perp PD\). 于是 \(PA\perp PG\). 同理，\(CD\perp PD\)，结合 \(PG\parallel CD\) 得 \(PD\perp PG\). 因此两平面的平面角可取 \(\angle APD\)，而 \(PA\perp PD\)，所以平面 \(PAB\perp\) 平面 \(PCD\).
\item 取 \(PC\) 的中点 \(H\)，连接 \(FH\)，\(DH\). 在三角形 \(PBC\) 中，\(F\)，\(H\) 分别是 \(PB\)，\(PC\) 的中点，所以 \(FH\parallel BC\)，且 \(FH=\dfrac12BC\). 在矩形 \(ABCD\) 中，\(DE\parallel BC\)，且 \(DE=\dfrac12AD=\dfrac12BC\). 因此 \(DE\parallel FH\) 且 \(DE=FH\)，四边形 \(EFHD\) 是平行四边形，从而 \(EF\parallel DH\). 又 \(DH\) 在平面 \(PCD\) 内，而 \(EF\) 不在该平面内，故 \(EF\parallel\) 平面 \(PCD\).
\end{enumerate}
\end{solution}


\begin{problem}

设函数 \(f(x)=\left[ax^2-(3a+1)x+3a+2\right]\e^x\).

\begin{enumerate}
\item 若曲线 \(y=f(x)\) 在点 \(\left(2,f(2)\right)\) 处的切线斜率为 \(0\)，求 \(a\)；
\item 若 \(f(x)\) 在 \(x=1\) 处取得极小值，求 \(a\) 的取值范围.
\end{enumerate}
\end{problem}

\begin{solution}
令 \(g(x)=ax^2-(3a+1)x+3a+2\)，则 \(f(x)=g(x)\e^x\)，所以
\[
f'(x)=\bigl(g'(x)+g(x)\bigr)\e^x=\bigl[ax^2-(a+1)x+1\bigr]\e^x=(x-1)(ax-1)\e^x
\]

\begin{enumerate}
\item 切线斜率为 \(f'(2)=0\)，故 \((2-1)(2a-1)\e^2=0\)，从而 \(a=\dfrac12\).
\item 因为 \(\e^x>0\)，只需讨论 \((x-1)(ax-1)\) 的符号. 若 \(a=0\)，则 \(x=1\) 左侧导数为正、右侧为负，取得极大值；若 \(a=1\)，则 \(f'(x)=(x-1)^2\e^x\geq0\)，没有极值.

当 \(a>1\) 时，\(\dfrac1a<1\)，在 \(\left(\dfrac1a,1\right)\) 上 \(f'(x)<0\)，在 \((1,+\infty)\) 上 \(f'(x)>0\)，所以 \(x=1\) 处取得极小值. 当 \(0<a<1\) 时，\(\dfrac1a>1\)，在 \((1,\dfrac1a)\) 上导数为负，故 \(x=1\) 处取得极大值. 当 \(a<0\) 时，两个零点为 \(\dfrac1a<1\) 和 \(1\). 在 \(\left(-\infty,\dfrac1a\right)\)、\(\left(\dfrac1a,1\right)\)、\((1,+\infty)\) 上，导数的符号依次为负、正、负，所以 \(x=1\) 处取得极大值. 因此 \(a\) 的取值范围为 \((1,+\infty)\).
\end{enumerate}
\end{solution}


\begin{problem}

已知椭圆 \(M:\dfrac{x^2}{a^2}+\dfrac{y^2}{b^2}=1\)（\(a>b>0\)）的离心率为 \(\dfrac{\sqrt6}{3}\)，焦距为 \(2\sqrt2\). 斜率为 \(k\) 的直线 \(l\) 与椭圆 \(M\) 有两个不同的交点 \(A\)，\(B\).

\begin{enumerate}
\item 求椭圆 \(M\) 的方程；
\item 若 \(k=1\)，求 \(\lvert AB\rvert\) 的最大值；
\item 设 \(P(-2,0)\)，直线 \(PA\) 与椭圆 \(M\) 的另一个交点为 \(C\)，直线 \(PB\) 与椭圆 \(M\) 的另一个交点为 \(D\). 若 \(C\)，\(D\) 和点 \(Q\left(-\dfrac74,\dfrac14\right)\) 共线，求 \(k\).
\end{enumerate}
\end{problem}

\begin{solution}
\begin{enumerate}
\item 设椭圆的半焦距为 \(c\). 由焦距为 \(2\sqrt2\)，得 \(c=\sqrt2\). 又由离心率
\[
\frac ca=\frac{\sqrt6}{3}
\]
得
\[
a=\frac{c}{\sqrt6/3}=\sqrt3
\]
因此
\[
b^2=a^2-c^2=3-2=1
\]
故椭圆 \(M\) 的方程为
\[
\frac{x^2}{3}+y^2=1
\]

\item 当 \(k=1\) 时，设直线 \(l\) 的方程为 \(y=x+m\). 将其与椭圆方程联立，得
\[
\frac{x^2}{3}+(x+m)^2=1
\]
即
\[
4x^2+6mx+3m^2-3=0
\]
设交点 \(A\)，\(B\) 的横坐标为 \(x_1\)，\(x_2\)，则
\(x_1+x_2=-\frac{3m}{2}\)，\(x_1x_2=\frac{3m^2-3}{4}\).
因为 \(A\)，\(B\) 为两个不同的交点，判别式应满足
\[
(6m)^2-4\cdot4(3m^2-3)=48-12m^2>0
\]
所以 \(|m|<2\). 又因直线斜率为 \(1\)，
\[
|AB|=\sqrt{(x_1-x_2)^2+\bigl((x_1+m)-(x_2+m)\bigr)^2}=\sqrt2\,|x_1-x_2|
\]
由根的判别式，
\[
(x_1-x_2)^2=(x_1+x_2)^2-4x_1x_2
=\frac{9m^2}{4}-3(m^2-1)=3-\frac{3m^2}{4}\leq3
\]
因此 \(|AB|\leq\sqrt6\)，当且仅当 \(m=0\) 时取等号，故 \(|AB|\) 的最大值为 \(\sqrt6\).

\item 设椭圆上一点为 \(X(x_0,y_0)\). 过点 \(P(-2,0)\) 和 \(X\) 的直线方程为
\[
y=\frac{y_0}{x_0+2}(x+2)
\]
这里 \(x_0\in[-\sqrt3,\sqrt3]\)，所以 \(x_0+2>0\)，且 \(4x_0+7\geq7-4\sqrt3>0\). 将直线方程代入椭圆方程，并利用 \(X\) 在椭圆上满足 \(x_0^2+3y_0^2=3\)，得到交点横坐标 \(x\) 所满足的方程
\[
(4x_0+7)x^2+12y_0^2x+12y_0^2-3(x_0+2)^2=0
\]
化简常数项，并用韦达定理计算两根之和，得
\(x+x_0=-\frac{12y_0^2}{4x_0+7}\)，\(x_0^2+3y_0^2=3\)
从而另一个交点 \(T\) 的坐标为
\(x_T=-\frac{7x_0+12}{4x_0+7}\)，\(y_T=\frac{y_0}{4x_0+7}\).
因此，若 \(A=(x_1,y_1)\)，\(B=(x_2,y_2)\)，则
\(C=\left(-\frac{7x_1+12}{4x_1+7},\frac{y_1}{4x_1+7}\right)\)，\(D=\left(-\frac{7x_2+12}{4x_2+7},\frac{y_2}{4x_2+7}\right)\).
由 \(Q\left(-\dfrac74,\dfrac14\right)\) 的坐标，计算得
\[
\overrightarrow{QC}=\frac1{4(4x_1+7)}(1,\,4y_1-4x_1-7)
\]
\[
\overrightarrow{QD}=\frac1{4(4x_2+7)}(1,\,4y_2-4x_2-7)
\]
由于 \(C\)，\(D\)，\(Q\) 共线，且上述向量的第一坐标均非零，所以
\[
4y_1-4x_1-7=4y_2-4x_2-7
\]
即
\[
y_1-x_1=y_2-x_2
\]
又 \(A\)，\(B\) 是不同的交点，故 \(x_1\ne x_2\)，从而直线 \(AB\) 的斜率为
\[
k=\frac{y_1-y_2}{x_1-x_2}=1
\]
\end{enumerate}
\end{solution}
